% Synheart HSI -- Kinematic Axis: Intern Progress Report (<= 3 pages)
% Compile:  pdflatex progress_report
\documentclass[11pt]{article}

\usepackage[utf8]{inputenc}
\usepackage[T1]{fontenc}
\usepackage[margin=0.85in]{geometry}
\usepackage{amsmath}
\usepackage{amssymb}
\usepackage{graphicx}
\usepackage{enumitem}
\usepackage[hidelinks]{hyperref}

\setlist[itemize]{leftmargin=1.4em, itemsep=2pt, topsep=2pt, parsep=0pt}
\setlength{\parindent}{0pt}
\setlength{\parskip}{0.4em}
\setlength{\emergencystretch}{2.5em}
\newcommand{\nb}[1]{\texttt{\footnotesize #1}}

% one validation panel: image + a short result line underneath
\newcommand{\panel}[2]{%
  \begin{minipage}[t]{0.48\textwidth}\centering
    \includegraphics[width=\linewidth]{#1}\\[2pt]
    {\footnotesize #2}
  \end{minipage}}

\pagestyle{plain}

\begin{document}

\begin{center}
  {\large\bfseries Synheart HSI --- Kinematic Axis: Intern Progress Report}\\[3pt]
  % TODO: confirm author name
  Bemnet Mussa \quad\textbar\quad Synheart AI \quad\textbar\quad \today
\end{center}
\vspace{-0.3em}\hrule\vspace{0.6em}

% ----------------------------------------------------------------------------
\section*{1.\quad Task}

The kinematic axis of Synheart's Human State Interface (HSI): infer a person's physical
state from a phone's motion and field sensors, under three constraints --- \emph{on-device},
\emph{deterministic} (rule-based, no machine learning), and \emph{privacy-first} (only
state readings leave the device, never the raw signal). I was assigned the Activity State
and Movement Regularity axes, and extended the work to the full set of five deterministic
states plus a single live system that runs all of them together off one phone stream.

% ----------------------------------------------------------------------------
\section*{2.\quad The five states and how they validate}

Each state is a deterministic rule over an explicit physical quantity, built and validated
on public datasets (MotionSense, HHAR, PAMAP2). The panels below show, for each state, the
plot from its notebook that demonstrates the rule works, with the measured result. Where a
part is not yet tested, it is called out honestly.

\vspace{0.4em}
\panel{figs/act_features.png}{\textbf{activity\_state --- 97.1\%} on MotionSense
(sed/walk/run; macro-F1 0.95, ECE 0.024), \textbf{99.6\%} cross-dataset on HHAR, and
\textbf{93.7\%} once the cycling branch is added (sed/walk/cycling, HHAR). The magnitude
thresholds separate the classes orientation-invariantly; cycling is caught by gyroscope
energy in the low-motion region.}
\hfill
\panel{figs/reg_dist.png}{\textbf{movement\_regularity --- separation, not accuracy.} A
continuous score from the Moe-Nilssen trunk-accelerometry autocorrelation method
(\emph{J.\ Biomechanics} 2004). No labelled ground truth, so no accuracy by design;
validated by clean separation (rhythmic gait $\approx$0.9 vs static $\approx$0.1) and
ROC-AUC 0.95 rhythmic-vs-irregular on PAMAP2.}

\vspace{0.8em}
\panel{figs/pos_acc.png}{\textbf{postural\_state --- 95\%} (lying / upright / in-motion,
PAMAP2 chest). The one axis that keeps orientation: posture \emph{is} gravity direction, so it
measures the tilt between current gravity and a calibrated upright. Accurate because lying
swings gravity $\sim$85$^\circ$ off upright (a wide, clean gap) and the threshold sits above
the whole upright spread, so upright is never misread as lying --- $\sim$97\% on the core
lying-vs-upright question, $\sim$94\% on a leak-proof held-out test. Chest + calibration is an
upper bound; the pocket deployment splits sitting/standing at the thigh.}
\hfill
\panel{figs/loco_pair.png}{\textbf{locomotion\_state --- to be tested.} The live rule reads
step rhythm for \emph{on foot} and a self-calibrated magnetometer for the hard
stationary-vs-in-vehicle split --- a steel cabin distorts the geomagnetic field, giving a
clean field-magnitude / inclination gap (shown, single-car PAMAP2 demo). This is a
demonstration, not a validation; proper testing needs a transport-mode dataset (SHL).}

\vspace{0.8em}
\begin{center}
\panel{figs/overall.png}{\textbf{overall\_state --- the combiner.} A deterministic priority
cascade turns the four states into a situation + exertion (0--3) + heart-rate note. The
notebook demonstrates the logic end-to-end on PAMAP2 at \textbf{71\%} (with simplified
stand-in axes; errors inherit from locomotion). The \emph{live} combiner runs the real four
axes and adds a cycling situation, so it is not identical to this 71\% demo.}
\end{center}

\vspace{0.4em}
Holding it together, beyond the individual states:

\begin{itemize}
  \item \textbf{Prior-art grounding.} Each state maps to an established measurement method
        (acceleration energy / ENMO, Moe-Nilssen autocorrelation, tilt / magnetic
        inclination, magnitude orientation-invariance); the contribution is their
        deterministic integration, not a new model.
  \item \textbf{Convergence testing.} Running the five together off one pocket stream
        surfaced and fixed real bugs --- an $\mathrm{m/s^2}\!\to\! g$ unit mismatch, a
        thigh-tilt inversion, and magnetometer baselines that had to become self-calibrated
        at startup.
  \item \textbf{Live on-device system.} One deterministic $\sim$50~Hz pipeline runs all five
        states from a single phone stream (Phyphox over WiFi), with a terminal technical
        view and a human-facing web view, startup calibration (with a PASS/FAIL check), and
        graceful degradation when a sensor is missing. \emph{It has been run successfully on
        a real phone.}
  \item \textbf{Technical report.} A separate report documents the live implementation and
        the design approach.
\end{itemize}

% ----------------------------------------------------------------------------
\section*{3.\quad What is left / next}

\begin{itemize}
  \item \textbf{Per-scenario testing.} The live system runs on a physical phone; what
        remains is systematically testing each scenario on-device (each activity, posture,
        and locomotion mode), rather than only on recorded data.
  \item \textbf{Locomotion across vehicles.} Replace the single-vehicle PAMAP2 result with a
        proper multi-vehicle evaluation (SHL) and reduce the remaining stationary-vs-vehicle
        confusion.
  \item \textbf{Posture \nb{device\_resting} branch.} The flat-pocket distinction (lying vs
        a phone resting on a surface) is untested and currently flagged low-confidence.
  \item \textbf{Transition windows.} 5\,s windows that straddle an activity change mix two
        states; this is not yet characterized.
  \item \textbf{Provisional thresholds.} Several thresholds were carried from the dataset
        they were tuned on and need re-tuning on deployment (pocket) data.
  \item \textbf{Activity ``vigorous'' class.} Currently unmodeled.
\end{itemize}

\vspace{1em}
\noindent\rule{\linewidth}{0.4pt}\\[3pt]
{\footnotesize \textbf{Source.} Every number and plot in this report is reproduced from the
project notebooks (\nb{maths/}) and the live system (\nb{live/}); each state's notebook carries
its full derivation, plots, and validation. Public repository:
\url{https://github.com/BemnetMussa/Kinematics-axis}.}

\end{document}
